\documentclass[aspectratio=169,11pt]{beamer}
\usepackage[utf8]{inputenc}
\usepackage{amsmath,amssymb}
\usepackage{booktabs}
\usepackage{graphicx}
\usepackage{xcolor}
\usepackage{textcomp}
\usepackage{microtype}

\graphicspath{%
  {figs/}%
  {figs_results/}%
  {../plots/}%
  {../remote_data/t2p23/}%
}

\usetheme{default}
\useinnertheme{rectangles}
\setbeamertemplate{navigation symbols}{}
\setbeamertemplate{footline}[frame number]

\definecolor{ymblue}{RGB}{31,73,125}
\definecolor{ymgreen}{RGB}{27,175,122}
\definecolor{ymred}{RGB}{178,34,52}

\setbeamercolor{structure}{fg=ymblue}
\setbeamercolor{title}{fg=ymblue}
\setbeamercolor{frametitle}{fg=white,bg=ymblue}
\setbeamercolor{block title}{bg=ymblue!15,fg=ymblue!90!black}
\setbeamercolor{block body}{bg=ymblue!4}
\setbeamercolor{block title example}{bg=ymgreen!20,fg=ymgreen!40!black}
\setbeamercolor{block body example}{bg=ymgreen!6}
\setbeamercolor{block title alerted}{bg=ymred!15,fg=ymred!70!black}
\setbeamercolor{block body alerted}{bg=ymred!5}
\setbeamerfont{block title}{size=\normalsize}
\setbeamerfont{normal text}{size=\footnotesize}
\setbeamertemplate{itemize items}[circle]
\setlength{\parskip}{2pt}

\newenvironment{physpicture}{\begin{exampleblock}{Physical picture}}{\end{exampleblock}}
\newenvironment{physmeaning}{\begin{exampleblock}{Physical meaning}}{\end{exampleblock}}

\newcommand{\eps}{\varepsilon_s}
\newcommand{\Vo}{V_0}
\newcommand{\kz}{k_z}
\newcommand{\OmA}{\Omega_A}
\newcommand{\OmF}{\Omega_F}
\newcommand{\Az}{A_{z1}}
\newcommand{\Ceil}{\Gamma}
\newcommand{\Kwn}{K}
\newcommand{\ii}{\mathrm{i}}
\newcommand{\Rey}{\operatorname{Re}}
\newcommand{\Imy}{\operatorname{Im}}
\newcommand{\sech}{\operatorname{sech}}
\newcommand{\dd}{\mathrm{d}}
\newcommand{\D}{\mathcal{D}}

\title{\textbf{The Non-Abelian Kelvin--Helmholtz Instability}}
\subtitle{A shear-driven SU(2) plasma instability: theory and GPU validation}
\author{SU(2) Yang--Mills KH simulation program}
\date{2026-07-22}

\begin{document}

\begin{frame}
\titlepage
\end{frame}

% ==============================================================================
\begin{frame}{Motivation}
Two counter-streaming, color-charged plasma beams with a shear flow layer between
them, coupled to a non-Abelian SU(2) gauge field.
\begin{itemize}
\item Classical (Abelian) Kelvin--Helmholtz instability is textbook physics for
neutral shear flows. Its non-Abelian analogue is not.
\item Relevant wherever colored plasma has shear: early-universe electroweak
plasma, quark--gluon plasma, and proposed dark non-Abelian sectors.
\item \textbf{No prior work} treats a genuinely shear-driven, non-Abelian
instability with a full wavenumber map. The closest results in the literature
are either two-stream-driven (Manuel--Mr\'owczy\'nski) or anisotropy-driven
chromo-Weibel modes (Mr\'owczy\'nski et al.) — not shear.
\end{itemize}
\begin{physpicture}
This talk derives the instability from first principles, finds an exact
reduction to a single ODE, extracts a growth-rate ceiling and a
wavelength-selection rule, and confronts all of it with GPU simulation data.
\end{physpicture}
\end{frame}

% ==============================================================================
\section{The system}

\begin{frame}{The simulated system}
For each color $a\in\{1,2,3\}$: fields $E_x^a,E_z^a,B_y^a,A_z^a$. For each beam
$s\in\{A,B\}$: fluid density $n_s$, momenta $p_{xs},p_{zs}$, color charge $Q_s^a$.
\vspace{0.05em}
\begin{align*}
&\partial_t n+\nabla\!\cdot(n\vec v)=0, \qquad
\partial_t p_x+\nabla\!\cdot(\vec v p_x)=-n\textstyle\sum_a Q^a(E_x^a+v_zB_y^a)\\
&\partial_t p_z+\nabla\!\cdot(\vec v p_z)=-n\textstyle\sum_a Q^a(E_z^a-v_xB_y^a),\qquad
\partial_t Q^a=\alpha\,v_z(\vec Q\times\vec A_z)^a\\[0.3em]
&\partial_t E_x^a=-\partial_zB_y^a-\alpha(\vec A_z\times\vec B_y)^a-J_x^a,\qquad
\partial_t E_z^a=\partial_xB_y^a-J_z^a\\
&\partial_t B_y^a=\partial_xE_z^a-\partial_zE_x^a+\alpha(\vec A_z\times\vec E_x)^a,
\qquad \partial_t A_z^a=-E_z^a
\end{align*}
Beam A: $v_z=+\Vo\tanh\xi$, $Q^1=+1$; beam B: $v_z=-\Vo\tanh\xi$, $Q^1=-1$;
background potential $\Az(\xi)=-\Vo\ln\cosh\xi$ sets $u(\xi)\equiv-\alpha\Az(\xi)\ge0$.
$\alpha$ is the non-Abelian coupling; $\Vo$ the shear velocity; $\kz$ the
perturbation wavenumber along the flow.
\end{frame}

\begin{frame}{The six-field linear system}
\footnotesize
Linearizing about the sheared background in the color-2/3 sector (which carries
charge $\pm1$ under the residual color-1 axis) gives a closed system of six
complex fields, $b,e_x,e_z,a,q_A,q_B\propto e^{\ii\kz z+\gamma t}$:
\[
\boxed{
\begin{aligned}
\gamma\,b &= -\ii\,\OmF\,e_x+e_z' & \gamma\,a &= -e_z\\
\gamma\,e_x &= -\ii\,\OmA\,b & \gamma\,q_A &= \ii\alpha v\,a-\ii v\,\OmA\,q_A\\
\gamma\,e_z &= b'+v\,(q_A-q_B) & \gamma\,q_B &= \ii\alpha v\,a+\ii v\,\OmA\,q_B
\end{aligned}}
\qquad
\OmA=\kz-u(\xi),\;\; \OmF=\kz+u(\xi)
\]
\vspace{0.2em}
\textbf{Validated directly against the 2D simulation}: 1--4\% agreement on the
core validation series, 4--9\% median across the full wavenumber map.
\end{frame}

% ==============================================================================
\section{Two branches, one mechanism}

\begin{frame}{Two branches emerge from the same background}
Freezing the background locally, the system splits into an electromagnetic
block and a charge-precession block.
\[
\underbrace{\gamma^2=-\OmA\OmF=u^2-\kz^2}_{\text{outer, tachyonic branch}}
\qquad\qquad
\underbrace{\gamma^2=-g(\xi;\gamma),\quad g=\frac{2\alpha v^3\OmA}{\gamma^2+v^2\OmA^2}}_{\text{shear branch, inside the layer}}
\]
\centering\includegraphics[width=0.28\textwidth]{fig_tachyon.pdf}
\raggedright
\begin{physpicture}
\footnotesize
Far from the layer, the strong background alone is unstable — a
Nielsen--Olesen-type charged-vector mode, independent of the flow. Inside the
layer, the flow shakes the beams' color charge; their \emph{difference}
current feeds back into the field — a genuine shear-driven resonance. This
talk focuses on the shear branch.
\end{physpicture}
\end{frame}

\begin{frame}{The $\kz=0$ limit: a validated anchor}
\footnotesize
At $\kz=0$ the shear-branch dispersion relation reduces to a cubic with an
exact closed-form growth rate:
\[
\boxed{\;\gamma(\kz{=}0)=\Bigl(\sqrt{\tfrac{\alpha^3}{2}}\,\Vo\Bigr)^{1/3}\sin 60^\circ\;}
\]
\begin{physpicture}
Not a bunching instability — nothing varies along $z$. It is
\emph{filamentation}: the background field rotates color-2/3 fields into each
other while the beams' counter-streaming currents feed the growth. This is a
member of the well-known chromo-Weibel family, appearing here as the
$\kz\to0$ limit of the shear branch.
\end{physpicture}
\textbf{Validated to 0.5\%}: at $\alpha=2,\Vo=0.1$, simulated
$\gamma=0.5039$ vs.\ predicted $0.5065$, over nine decades of growth — this is
the strongest, parameter-free anchor for everything that follows, and it holds
on two independent simulation configurations three weeks apart.
\end{frame}

% ==============================================================================
\section{An exact reduction, and a growth ceiling}

\begin{frame}{Reduction to a single ODE}
\footnotesize
The six-field system has only one spatial derivative chain; every other field
is algebraic in $a(x)$, the color-2/3 potential perturbation. Eliminating them
gives one second-order equation:
\[
\boxed{\;\Bigl(\frac{\gamma^2 a'}{\D}\Bigr)'=(\gamma^2+g)\,a,\qquad
\D(x)=\gamma^2+\kz^2-u(x)^2,\qquad
g(x)=\frac{2\alpha v^3\OmA}{\gamma^2+v^2\OmA^2}\;}
\]
\centering\includegraphics[width=0.42\textwidth]{fig_anatomy.pdf}
\raggedright
\begin{physpicture}
$\D$ changes sign exactly where the two branches meet: $\D>0$ is a trapped
shear well, $\D<0$ is the tachyonic outer region — one mechanism, two fates.
$g$ is the only drive of the shear branch, and it dies as $v^3$ toward the
layer center: the wave must shake the charges, make the two beams respond
\emph{differently}, and collect that difference as a current.
\end{physpicture}
\textbf{Exactness verified}: the reduced equation reproduces the full
six-field eigenvalues to better than $10^{-10}$ relative difference.
\end{frame}

\begin{frame}{A growth-rate ceiling}
The drive $g$ is a resonance in the beam-frame precession frequency, maximized
when that frequency matches the growth rate itself. Its peak value is
$G_{\max}=\alpha\Vo^2/\gamma$. A confining well needs $G_{\max}>\gamma^2$
somewhere, which forces
\[
\boxed{\;\gamma^3\le\alpha\Vo^2\;}
\]
independent of $\kz$ and of the shear-layer width.
\begin{physpicture}
A self-limiting resonance: the beam response is a Lorentzian in local
precession frequency, and its width is set by $\gamma$ itself. The faster the
mode grows, the broader the band it averages over, and the weaker the peak
drive it can extract. The ceiling is the fixed point of that feedback.
\end{physpicture}
\textbf{Measured}: peak growth rates reach 95--99\% of $(\alpha\Vo^2)^{1/3}$ in
every simulated series tested (shown quantitatively later).
\end{frame}

\begin{frame}{A universal, dimensionless well}
\footnotesize
In natural units $\Ceil\equiv(\alpha\Vo^2)^{1/3}$ (rate scale) and
$\Kwn\equiv(\alpha/\Vo)^{1/3}$ (wavenumber scale), the drive becomes universal
— no $\alpha,\Vo$ remain:
\[
G(w)=\Ceil^2\,\hat G(\hat w),\qquad \hat G(\hat w)=\frac{2\hat w}{\hat\gamma^2+\hat w^2}
\]
\centering\includegraphics[width=0.44\textwidth]{fig_drive.pdf}
\raggedright
\textbf{Growth rates scale as $(\alpha\Vo^2)^{1/3}$ and wavenumbers as
$(\alpha/\Vo)^{1/3}$} — not as the naive product $\alpha\Vo$. This scaling is
directly confirmed on GPU data later in the talk.
\end{frame}

% ==============================================================================
\section{Wavelength selection}

\begin{frame}{Quantization of the shear well}
\footnotesize
The well is shallow — less than a quarter wavelength of accumulated phase — so
the growth rate is set by a connection condition at the well's edge rather
than a simple cavity-mode count:
\[
\tan S=\frac{\kappa}{k_{\rm edge}},\qquad \kappa=\sqrt{\gamma^2+\kz^2},\qquad
S=\int\sqrt{Q}\,dx
\]
\centering\includegraphics[width=0.4\textwidth]{fig_quant.pdf}
\raggedright
\begin{physpicture}
A surface-mode-like state, not a cavity mode: growth rates hug the ceiling (a
small drop in $\gamma$ buys all the depth needed), while the wavelength at
which growth peaks is set by the geometry of the confining window rather than
by the shear profile itself.
\end{physpicture}
Solving this quantization condition exactly reproduces the independently
computed six-field eigenvalues to $\le2\%$ (median $\le0.2\%$) across a wide
survey of $(\alpha,\Vo,\kz)$.
\end{frame}

\begin{frame}{No intrinsic wavelength maximum}
In an unbounded domain, nothing stops the mode from following the resonance
outward as $\kz$ grows: the peak drive $G_{\max}=\alpha\Vo^2/\gamma$ does not
depend on $\kz$, so $\gamma(\kz)$ climbs monotonically toward the ceiling —
directly confirmed: $\gamma$ reaches 99.5\% of the ceiling by $\kz=10$ in an
unbounded test, with no turnover.
\begin{physpicture}
The peaked dispersion curve seen in the finite simulation domain is therefore
\textbf{selected by the box, not intrinsic to the flow} — the single most
counter-intuitive result of the theory, and the one most directly tested
against data next.
\end{physpicture}
\end{frame}

\begin{frame}{The confinement-set wavelength}
\footnotesize
The resonance location moves outward linearly with $\kz$; setting it equal to
the domain's confinement radius $\xi_w$ gives the wavenumber at which growth
peaks:
\[
\boxed{\;\kz^{\rm peak}\simeq\alpha\Vo(\xi_w-\ln2)+c\,(\alpha/\Vo)^{1/3},\qquad c\approx1\;}
\]
\centering\includegraphics[width=0.42\textwidth]{fig_peak.pdf}
\raggedright
\textbf{Both terms are independent of the shear-layer width} — coupling
strength and confinement radius jointly select the wavelength; the flow
profile itself does not directly enter. This is the sharpest and most novel
claim of the theory, and is tested directly on GPU data next.
\end{frame}

% ==============================================================================
\section{Confronting the theory with GPU data}

\begin{frame}{Non-monotonic dispersion, on real simulation data}
\footnotesize
\begin{columns}[T]
\begin{column}{0.5\textwidth}
\centering
\includegraphics[width=\textwidth]{fig_r_dispersion.pdf}\\
fixed $\Vo=0.05$, varying $\alpha$
\end{column}
\begin{column}{0.5\textwidth}
\centering
\includegraphics[width=\textwidth]{fig_r_dispersion_v0.pdf}\\
fixed $\alpha=1$, varying $\Vo$
\end{column}
\end{columns}
\vspace{0.3em}
Points are simulated $\gamma(\kz)$; curves are the $\sigma$-chased
eigensolver's theoretical prediction, drawn underneath. Every curve peaks at
an intermediate $\kz$, and the peak location climbs with $\alpha$ — from
$\kz^{\rm peak}=2$ at $\alpha=0.4$ to $5$ at $\alpha=3.5$, at fixed shear width.
Classical Kelvin--Helmholtz also peaks, but its peak is fixed by the shear
width alone; a peak that moves with the gauge coupling, at fixed flow, has no
classical analogue.
\end{frame}

\begin{frame}{The peak tracks confinement, not the shear width}
\footnotesize
\begin{columns}[T]
\begin{column}{0.4\textwidth}
\includegraphics[width=\textwidth]{fig_r_kzpeak.pdf}
\end{column}
\begin{column}{0.58\textwidth}
\begin{tabular}{lcc}
$(\alpha,\Vo,\xi_w)$ & measured & formula\\\midrule
(1, 0.05, 20) & 3.5 & 3.7\\
(2, 0.05, 11) & 4.5 & 4.5\\
(3, 0.03, 5) & 5.5 & 5.0\\
(5, 0.03, 5) & 6.5 & 6.2\\
(1, 0.20, 10) & 3.0 & 3.6\\
\end{tabular}\\[0.3em]
\textbf{Direct window test}, same $(\alpha,\Vo)$: widening the confinement
radius alone moves the measured peak from $\kz=4.5$ to $5.5$, matching the
formula's $4.45\to5.85$.
\end{column}
\end{columns}
\begin{physpicture}
Two different confinement radii, same coupling and shear velocity, two
different — and both correctly predicted — peak wavenumbers: confinement and
coupling jointly select the wavelength, not the flow profile.
\end{physpicture}
\end{frame}

\begin{frame}{The dimensionless collapse}
\footnotesize
\centering
\includegraphics[width=0.72\textwidth]{t2p5_collapse.png}
\[
\gamma_{\rm peak}/(\alpha\Vo^2)^{1/3}=0.977\pm0.011\quad(1.1\%\text{ spread, }10\times\text{ range in }\alpha\Vo)
\]
\raggedright
The single most quantitative confirmation: fifteen independent dispersion
curves, spanning a factor of ten in $\alpha\Vo$, collapse onto one master
curve — direct proof that $\alpha,\Vo$ enter growth rates through
$(\alpha\Vo^2)^{1/3}$, exactly as the ceiling predicts.
\end{frame}

\begin{frame}{Full-map accuracy}
\footnotesize
Growth rates measured across a broad grid in $\alpha$, $\Vo$, and $\kz$
(1087/1134 points, 96\% coverage) against the independently computed
eigenvalue, at fixed shear width:

\begin{center}
\includegraphics[width=0.32\textwidth]{fig_r_error_by_v0.pdf}
\end{center}

Median relative error \textbf{4--8\%} depending on $\Vo$, improving
monotonically with higher shear velocity (int-tier points, $\alpha\in[0.4,1.0]$
common to every $\Vo$ so the bars are an apples-to-apples comparison rather
than an artifact of which $\alpha$'s each campaign happened to sweep). The
theory and the simulation agree across the whole mapped parameter space, not
just at isolated anchor points.
\end{frame}

\begin{frame}{Independence from the shear-layer width}
\begin{center}
\includegraphics[width=0.85\textwidth]{eps_collapse_gpu_validation.png}
\end{center}

\small
The full ceiling-and-peak theory, re-expressed in terms of a single
dimensionless group combining coupling, shear velocity, and layer width,
collapses 91 simulation runs spanning a wide range of shear-layer widths onto
one universal curve: median relative error 9.9\%. The wavelength-selection law
does not depend on the microscopic shear width — only on coupling and
confinement.
\end{frame}

\begin{frame}{The eigenfunction structure matches, not just the rate}
\centering
\includegraphics[width=0.68\textwidth]{t2p2_eigenfunction_overlay.png}
\small
Spatial profile of the dominant field, simulation vs.\ theory: correlation
$=1.0000$. A companion field, seeded at less than 1\% of the driving field's
amplitude and grown entirely through the instability's internal feedback
chain, independently reproduces the theory's non-trivial double-lobe shape —
an emergent structural match, not just a matched growth rate.
\end{frame}

% ==============================================================================
\section{Summary}

\begin{frame}{Where this leaves the physics}
\small
\begin{tabular}{p{0.32\textwidth}p{0.6\textwidth}}
\toprule
Result & Status\\\midrule
$\kz=0$ chromo-Weibel rate & Closed form, validated to 0.5\%\\
Six-field linear system & Validated to 1--9\% against simulation\\
Exact scalar reduction & Reproduces linear theory to $<10^{-10}$\\
Growth ceiling $\gamma^3\le\alpha\Vo^2$ & Measured at 95--99\% of ceiling; $0.977\pm0.011$ collapse\\
Quantization / dispersion curve & Matches theory to $\le2\%$ (median $\le0.2\%$)\\
Confinement-set wavelength selection & Confirmed directly by window tests and full parameter map\\
\bottomrule
\end{tabular}
\vspace{0.3em}
\centering\textbf{A shear-driven non-Abelian instability, derived exactly and
confirmed on over a thousand GPU simulation runs: coupling and confinement
jointly select the wavelength; the flow profile does not enter directly.}
\end{frame}

\begin{frame}{Context and next steps}
\begin{itemize}
\item No prior literature treats a genuinely shear-driven non-Abelian
instability with a full wavenumber map — this appears to be new physics, not
a variant of a known result.
\item Closest related work: chromo-Weibel instabilities (anisotropy-driven,
not shear); the outer tachyonic branch is a Nielsen--Olesen-type mode;
Abelian shear instabilities (Das \& Kaw; Csernai et al.\ in heavy-ion
collisions) are the closest classical analogue but lack the non-Abelian
wavelength-selection mechanism.
\item Remaining work: a citation-trail audit before submission, and a
dedicated follow-up study of the high-coupling regime, which is currently
kept outside the headline claims while its numerical behavior is
characterized more fully.
\end{itemize}
\vspace{0.2em}
\centering\Large\textbf{The theory is exact, matches simulation to a few
percent, and is confirmed across a broad parameter survey.}
\end{frame}

\begin{frame}
\centering
\Huge Thank you.
\vspace{0.15em}

\normalsize\normalfont Questions, and where to push next.
\end{frame}

\end{document}
